Species interactions structure ecological communities, but predicting when interactions occur, how species are positioned within interaction networks, and how interactions scale up to community-level outcomes remain central challenges in ecology. The architecture of species interactions emerges from processes operating across multiple scales, from species-level traits to community structure to biogeographic contexts. In this dissertation, I examine how interspecific interactions are organized across these scales in mutualistic, antagonistic, and experimental systems. At the species level, I show that low-dimensional representations of species positions within avian frugivory networks can be used in concert with species traits and phylogenetic information to improve predictions of missing links. I also show that, for some taxa, both the geographic range size and the climatic niche breadth of a species are associated with its role in local and regional interaction webs. At broader scales, I use a global set of insular host-helminth interaction networks to show how island characteristics like isolation and area are associated with differences in species assemblages and network structure. Additionally, I review the current state of knowledge on how major classes of anthropogenic impacts affect the structural properties of plant-animal transport networks globally. Finally, I show how interspecific interactions scale to community-level outcomes under alternative environmental regimes in experimental microbial communities. By linking predictive models, network analyses, spatial approaches, synthesis, and experiments, this dissertation contributes to a multiscale view of how ecological interactions are organized across systems.
